% Part: incompleteness
% Chapter: theories-computability
% Section: oconsis-ext-of-q-undec

\documentclass[../../../include/open-logic-section]{subfiles}

\begin{document}

\olfileid{inc}{tcp}{oqn}

\olsection{$\Th{Q}$નાં $\omega$-સુસંગત વિસ્તરણો અનિર્ણેય છે}

\begin{explain}
$\Th{Q}$ સંગણકીય રીતે
પરિગણનીય ગણોમાં પૂર્ણ છે
તેની સાબિતી આ તથ્ય પર
આધાર રાખતી હતી કે
$\Th{Q}$માં સાબિત થતું
દરેક વાક્ય પ્રાકૃતિક
સંખ્યાઓ વિશે ``સાચું'' છે.
આગળની વ્યાખ્યા અને પ્રમેય
ઉપરની સાબિતી માટે જરૂરી
``સત્ય''ના ચોક્કસ પાસાઓ
ઓળખીને તે પરિણામને વધુ
મજબૂત બનાવે છે. જોકે આ
પ્રમેય પર વધારે સમય રોકાવાની
જરૂર નથી, કારણ કે ટૂંકમાં
તેને હજુ વધુ મજબૂત કરીશું.
અહીં તેને મુખ્યત્વે ઐતિહાસિક
હેતુથી સામેલ કર્યો છે:
ગડેલના મૂળ લેખમાં
$\omega$-સુસંગતતાની
વિભાવના વપરાઈ હતી, પરંતુ
થોડા સમય પછી સામાન્ય
સુસંગતતા વડે
$\omega$-સુસંગતતાને
બદલીને તેનું પરિણામ વધુ
મજબૂત બનાવાયું.
\end{explain}

\begin{defn}
\ollabel{thm:oconsis-q}
સિદ્ધાંત~$\Th{T}$ને
$\omega$-સુસંગત ત્યારે
કહીએ જ્યારે નીચેની શરત હોય:
$\lexists[x][!A(x)]$ કોઈપણ
વાક્ય હોય અને $\Th{T}$માં
$\lnot !A(\num 0)$,
$\lnot !A(\num 1)$,
$\lnot !A(\num 2)$,
\dots સાબિત થાય, તો
$\Th{T}$માં
$\lexists[x][!A(x)]$
સાબિત ન થાય.
\end{defn}

\begin{thm}
  $\Th{T}$ એ $\Th{Q}$ને
  સમાવતા કોઈપણ
  $\omega$-સુસંગત સિદ્ધાંત
  હોય. તો $\Th{T}$
  નિર્ણેય નથી.
\end{thm}

\begin{proof}
$\Th{T}$ એ $\Th{Q}$ને
સમાવતું હોય, તો
$\Th{T}$ સંગણનીય વિધેયો
અને સંબંધોને નિરૂપે છે.
અગાઉની સાબિતીમાં માત્ર
થોડો ફેરફાર કરવો પડે.
ઉપરની જેમ $x \in K$ હોય,
તો $\Th{T}$માં
$\lexists[s][!A_T(\num
  x,\num x, s)]$ સાબિત થાય.
વિપરીત રીતે માનો કે
$\Th{T}$માં
$\lexists[s][!A_T(\num x, \num x, s)]$
સાબિત થાય છે. તો $x$ એ
$K$માં હોવું જ પડે: અન્યથા
$x$ મશીનની $x$ ઇનપુટ
પર અટકતી સંગણના જ ન હોય.
$!A_T$ ક્લીનીના~$T$
સંબંધને નિરૂપે છે, એટલે
$\Th{T}$માં
$\lnot !A_T(\num x, \num x, \num 0)$,
$\lnot !A_T(\num x, \num
x, \num 1)$, \dots સાબિત
થાય. આથી $\Th{T}$
$\omega$-અસુસંગત થઈ જાય,
જે ધારણાની વિરુદ્ધ છે.
\end{proof}

\end{document}
